% !TEX root = ../../../main.tex
% 来源：中学数学实验教材·第一册（代数）
% 章节：一次方程式 / 解应用问题
% 原始文件：2.tex，行 2756-2800
% 类型：tikz
% ---
\begin{tikzpicture}[>=latex,scale=.7]
\begin{scope}
\draw[pattern=north east lines] (2,1.5) circle (2.5);
\draw[fill=white] (0,0) rectangle (4,3);
\draw[<->](0,0)--node[above]{$2R$}(4,3);
\node at (2,0)[above]{$a$};
\node at (0,1.5)[right]{$b$};

\end{scope}    
\begin{scope}[xshift=6cm]
    \draw[pattern=north east lines] (0,0) rectangle (5,3);
    \draw[fill=white]  (.5,0) arc (180:0:2);
    \draw [->] (2.5,0)--node[above]{$R$}+(40:2);
    \node at (2.5,3)[above]{$a$};
    \node at (5,1.5)[right]{$b$};
\end{scope} 
\begin{scope}[yshift=-6cm]
    \draw[pattern=north east lines]  (0,0) arc (180:0:2.6); 
    \draw[fill=white](2.6,1.3) circle(1.3);
    \draw (0,0)--(5.2,0);
    \draw[|<->|] (0,-.2)--node[below]{$D$}(5.2,-.2);
\end{scope} 
\begin{scope}[xshift=6cm, yshift=-6cm]

\fill  [pattern=north east lines] (0,0) rectangle (6,4);
\fill [white](0,0) rectangle (1,1);
\fill [white](0,3) rectangle (1,4);
\fill [white](5,0) rectangle (6,1);
\fill [white](5,3) rectangle (6,4);

\foreach \x in {0,4}
{
    \draw (1,\x)--(5,\x);
    \draw (\x*1.5, 1)--(\x*1.5, 3);
}

\draw (0,1)--(1,1)--node[left]{$x$}(1,0);
\draw (0,3)--(1,3)--node[left]{$x$}(1,4);
\draw (5,0)--(5,1)--(6,1);
\draw (5,4)--(5,3)--node[above]{$x$}(6,3);
\node at (3,4)[above]{$4x$};
\node at (6,2)[right]{$2x$};

\end{scope}     
\end{tikzpicture}
